% P17-D2: Rubinstein's finite-prize expected-utility presentation.
\chapter{Finite Lotteries and Expected Utility}
\label{sec:p17-finite-expected-utility}

A choice under risk compares probability distributions of consequences.
To evaluate such a distribution by averaging consequence values requires
more structure than an ordering of the consequences alone. This chapter
gives the lottery domain, the mixture and continuity assumptions, and the
complete representation argument for the finite-prize presentation in
Rubinstein's Lecture~7 \cite[2015 revision, pp.~94--100]{Rubinstein2015LectureNotes}.
It then turns the representation into a choice procedure and shows what
fails when an assumption is removed.

The probabilities in this model are supplied as part of each alternative.
Their estimation and the justification of a person's lottery preferences
are additional questions. The probability calculations use the finite
mass and expectation definitions of \cref{sec:p15-finite-expectation};
the ordering of certain prizes uses \cref{sec:p17-finite-preferences}.
The space of lotteries will generally be infinite even though the set of
prizes is finite.

\section{Consequences, distributions and mixtures}
\label{sec:p17-lottery-domain}

\begin{hbdefinition}{def:p17-finite-lottery}{A lottery on a finite prize set}
Let \(Z\) be a finite nonempty set of consequences, called \emph{prizes}.
Put
\[
 \Delta(Z)\defeq\left\{p:Z\to[0,1]:\sum_{z\in Z}p(z)=1\right\}.
\]
A member \(p\) is a \emph{lottery}; \(p(z)\) is the supplied probability
of prize \(z\). For \(z\in Z\), the \emph{certain lottery} \(\delta_z\)
has value one at \(z\) and zero at every other prize. The notation replaces
the source's \([z]\), leaving the same object. The source writes
\(L(Z)\) for this whole domain \cite[p.~94]{Rubinstein2015LectureNotes}.

For \(p,q\in\Delta(Z)\) and \(\alpha\in[0,1]\), define the
\emph{mixture} coordinate by coordinate:
\[
 [\alpha p+(1-\alpha)q](z)\defeq\alpha p(z)+(1-\alpha)q(z).
\]
More generally, for \(K\geq1\), lotteries \(p_1,\ldots,p_K\) and
weights \(\alpha_k\geq0\) with \(\sum_{k=1}^K\alpha_k=1\), write
\(\sum_{k=1}^K\alpha_kp_k\) for their coordinatewise mixture.
Zero-weight branches are permitted. An empty family cannot have weights
summing to one and is not a lottery mixture under this definition.
The source uses the mixture sign \(\oplus\); the addition sign here
denotes the same vector operation \cite[pp.~96--97]{Rubinstein2015LectureNotes}.
\index{lottery!finite prize set}\index{lottery!mixture}
\end{hbdefinition}

\begin{hbtheorem}{thm:p17-lottery-reduction}{proved=annex:p17.lottery-reduction}{Mixture probabilities and reduction}
Every defined mixture is in \(\Delta(Z)\). Every lottery has the
decomposition
\[
                         p=\sum_{z\in Z}p(z)\delta_z.
\]
For each \(k\), let \(\beta_{k\ell}\geq0\),
\(\sum_{\ell=1}^{J_k}\beta_{k\ell}=1\), and
\(p_{k\ell}\in\Delta(Z)\), with \(J_k\geq1\). Then
\[
 \sum_{k=1}^K\alpha_k
       \left(\sum_{\ell=1}^{J_k}\beta_{k\ell}p_{k\ell}\right)
 =\sum_{k=1}^K\sum_{\ell=1}^{J_k}
                       \alpha_k\beta_{k\ell}p_{k\ell}.
\]
In particular, a two-stage experiment with branch probabilities
\(\alpha_k\) and conditional prize laws \(p_k\) has final prize law
\(\sum_k\alpha_kp_k\). Its joint branch/prize mass is
\(\alpha_kp_k(z)\). A zero-probability branch does not affect that law.
\end{hbtheorem}

This reduction is an identity of probability laws. Treating two procedures
with the same final law as the same alternative is a modeling assumption
embodied in the domain \(\Delta(Z)\). Rubinstein's construction describes
separate randomization stages \cite[pp.~96--97]{Rubinstein2015LectureNotes}.
The conditional-law formulation above does not assert that the selected
branch and the final prize are independent: their joint mass generally
does not factor into their two marginal masses.

\begin{hbwarning}{def:p17-lottery-description-boundary}{A consequence must include what matters to the comparison}
Two actions can have the same distribution of a coarse outcome and still
produce different relevant consequences. For example, receiving a tool
when it is needed differs from receiving it only when it is unnecessary,
even when the unconditional probability of receiving it is the same.
The prize description must include the relevant circumstances before a
preference on its distribution can represent this comparison. This is the
issue illustrated by the source's rain-and-umbrella example
\cite[pp.~94--95]{Rubinstein2015LectureNotes}.

Likewise, a preference for the fairness of a random procedure can depend on
more than its realized prize. The lottery representation does not establish
that such a concern should be omitted. It states which preference domain
and mixture comparisons the following theorem assumes.
\end{hbwarning}

\section{The assumptions on lottery preferences}
\label{sec:p17-lottery-axioms}

\begin{hbdefinition}{def:p17-lottery-preference}{Weak preference on the lottery space}
A relation \(\succeq\) on \(\Delta(Z)\) is a \emph{weak preference}
if it is complete and transitive on that whole domain. Thus every pair
\(p,q\) has \(p\succeq q\) or \(q\succeq p\), including equal
arguments, and two successive weak comparisons imply the third.
Define \(p\succ q\) by \(p\succeq q\) and not \(q\succeq p\),
and \(p\sim q\) by both weak comparisons.

These are the relation conventions of \cref{def:p17-weak-preference}, now
applied to an infinite lottery domain when \(|Z|\geq2\). The earlier
finite ordinal-representation theorem alone does not produce a
representation on this domain. The additional axioms below are needed.
\end{hbdefinition}

\begin{hbdefinition}{def:p17-lottery-independence}{Independence of a common mixture component}
The \emph{independence axiom} requires, for all \(p,q,r\in\Delta(Z)\)
and \(0<\alpha<1\),
\[
 p\succeq q\quad\Longleftrightarrow\quad
 \alpha p+(1-\alpha)r\succeq\alpha q+(1-\alpha)r.
\]
This is the source's Axiom I \cite[p.~97]{Rubinstein2015LectureNotes}.
It compares alternatives with the same mixture weight and the same
remaining component. It is an axiom of preference, distinct from
probabilistic independence of events or random variables.

At \(\alpha=1\) the displayed comparison is unchanged directly.
At \(\alpha=0\) both mixtures equal \(r\), so the reverse implication
would incorrectly force every pair of lotteries to be indifferent.
That endpoint is deliberately excluded from the axiom.
\index{expected utility!independence axiom}
\end{hbdefinition}

\begin{hbdefinition}{def:p17-lottery-continuity}{Continuity in lottery probabilities}
Give \(\Delta(Z)\) its relative Euclidean distance
\[
 d(p,q)\defeq\left(\sum_{z\in Z}(p(z)-q(z))^2\right)^{1/2}.
\]
A weak preference is \emph{continuous} when, whenever \(p\succ q\),
there is an \(\varepsilon>0\) such that every pair of lotteries
\(p',q'\) with \(d(p',p)<\varepsilon\) and
\(d(q',q)<\varepsilon\) satisfies \(p'\succ q'\).
Only vectors that remain lotteries are quantified over; boundary points
such as certain lotteries have neighborhoods relative to \(\Delta(Z)\).
This is the source's neighborhood Axiom C, with one common radius chosen
inside its two neighborhoods \cite[p.~98]{Rubinstein2015LectureNotes}.
\index{expected utility!continuity axiom}
\end{hbdefinition}

\begin{hblemma}{thm:p17-mixture-comparison}{proved=annex:p17.mixture-comparison}{Replacement and ordering along a mixture segment}
Let \(\succeq\) be a weak preference satisfying independence.
Replacing a component of positive weight in a finite mixture by a weakly
preferred lottery weakly improves the mixture; strict improvement of that
component strictly improves it. An indifferent replacement preserves
indifference. Replacing a zero-weight component leaves the vector unchanged.
Consequently these comparisons can be applied successively to any finite
list of components.

If \(a\succ b\), put \(B_t\defeq ta+(1-t)b\) for \(0\leq t\leq1\).
For \(1\geq s>t\geq0\), one has \(B_s\succ B_t\).
If \(a\sim b\), every such \(B_t\) is indifferent to both endpoints.
The endpoint and zero-weight cases are included in these assertions.
\end{hblemma}

\begin{hblemma}{thm:p17-lottery-calibration}{proved=annex:p17.lottery-calibration}{Calibration between a better and a worse lottery}
Suppose the weak preference satisfies independence and continuity.
If \(a\succ b\) and \(a\succeq p\succeq b\), there is a unique
\(t\in[0,1]\) for which
\[
                            p\sim ta+(1-t)b.
\]
The value is zero exactly when \(p\sim b\), one exactly when
\(p\sim a\), and lies strictly between them when
\(a\succ p\succ b\). This proves the interpolation property called
\(C^*\) in the source and supplies the endpoint cases used in its
representation proof \cite[pp.~98--100]{Rubinstein2015LectureNotes}.
\end{hblemma}

\section{Representation and the permissible changes of scale}
\label{sec:p17-expected-utility-representation}

\begin{hbdefinition}{def:p17-expected-utility}{Consequence utility and expected utility of a lottery}
A function \(u:Z\to\mathbb R\) assigns a \emph{consequence utility}
to each prize. Its \emph{expected-utility function} is
\[
                U_u(p)\defeq\sum_{z\in Z}p(z)u(z)
                         \qquad(p\in\Delta(Z)).
\]
It represents \(\succeq\) when
\(p\succeq q\) holds exactly when \(U_u(p)\geq U_u(q)\).
The numbers \(u(z)=U_u(\delta_z)\) are also called von Neumann--Morgenstern
utilities or Bernoulli utilities. The source distinguishes its prize
function \(v\) from its lottery function \(U\)
\cite[p.~99]{Rubinstein2015LectureNotes}; here the former is denoted \(u\).
On the finite probability space \(Z\) with masses \(p(z)\),
\(U_u(p)\) is precisely \(\mathbb E_p[u]\).
\index{expected utility!finite representation}
\index{utility!von Neumann--Morgenstern}
\end{hbdefinition}

\begin{hbtheorem}{thm:p17-finite-vnm}{proved=annex:p17.finite-vnm}{Finite-prize expected-utility representation}
Let \(Z\) be finite and nonempty. A relation on \(\Delta(Z)\) is a
complete, transitive preference satisfying independence and continuity
if and only if some \(u:Z\to\mathbb R\) represents it by expected utility.

If all certain prizes are indifferent, all lotteries are indifferent and
one may take \(u(z)=0\) for every \(z\). This includes \(|Z|=1\).
Otherwise choose a best and a worst certain prize, \(z^+,z^-\), and
put \(M=\delta_{z^+}\), \(m=\delta_{z^-}\). They satisfy \(M\succ m\).
The representing function can be chosen uniquely with
\[
 u(z^+)=1,\qquad u(z^-)=0,\qquad
       \delta_z\sim u(z)M+(1-u(z))m\quad(z\in Z).
\]
For this choice, \(0\leq u(z)\leq1\), and every lottery satisfies
\[
                  p\sim U_u(p)M+(1-U_u(p))m.
\]
The uniqueness here fixes the chosen normalization and the entire
lottery preference, rather than merely an ordering of certain prizes.
\end{hbtheorem}

The proof follows the best/worst-prize construction of Rubinstein's vNM
theorem \cite[pp.~99--100]{Rubinstein2015LectureNotes}, with explicit
mixture replacement, interval calibration and all-indifferent cases.
It uses the completeness of the real interval in the calibration proof.
It does not require a separate general utility-representation theorem on
the lottery space.

\begin{hbtheorem}{thm:p17-vnm-affine-uniqueness}{proved=annex:p17.vnm-affine-uniqueness}{Positive affine uniqueness of consequence utilities}
Suppose \(u,v:Z\to\mathbb R\), with finite nonempty \(Z\), represent
the same lottery preference by expected utility. Then there are real
constants \(a>0\) and \(b\) such that
\[
                           v(z)=a u(z)+b\qquad(z\in Z).
\]
Conversely, every such transformation preserves all lottery comparisons
and ties, since \(U_v(p)=aU_u(p)+b\).
If \(u\) is nonconstant, the two constants are fixed by its values and
the corresponding \(v\)-values at a best and a worst prize. If \(u\)
is constant, \(v\) must be constant too, and \(a\) may be any positive
number with \(b\) chosen to match the two constants.
\end{hbtheorem}

A strictly increasing transformation \(f(U_u(p))\) of the \emph{whole lottery
score} preserves its ordering. In contrast, transforming the individual
prize values to \(f(u(z))\) and then averaging can change the ordering.
The positive affine theorem concerns this latter, consequence-level
representation \cite[p.~100]{Rubinstein2015LectureNotes}.

\section{Using the representation in a finite choice}
\label{sec:p17-expected-utility-use}

\begin{hbmethod}{def:p17-expected-utility-choice}{Evaluate an explicit menu of lotteries}
Given a finite prize set \(Z\), a finite menu \(\mathcal B\subseteq\Delta(Z)\),
and an expected-utility representation \(u\) of the intended preference:
\begin{enumerate}
\item Specify each consequence sufficiently to include the attributes being
compared. Record each lottery's probability mass and check normalization.
\item State the preference assumptions that license using \(u\). A ranking
of certain prizes alone does not determine its intermediate values.
Under a nontrivial model, a certain prize's indifference probability between
the fixed best and worst prizes supplies its normalized value.
\item Calculate \(U_u(p)=\sum_zp(z)u(z)\) for every menu member.
\item Retain every maximizer. Report ties; selecting one tied member needs
an additional choice rule. An empty menu has no feasible choice.
\item Examine how changing an uncertain probability or an intermediate
utility value changes the comparison. Distinguish the given model's
calculation from evidence that the model fits a decision maker.
\end{enumerate}
For a nonempty finite menu a maximum exists, and the representation theorem
identifies its maximizers exactly with the most preferred feasible lotteries.
This is the finite-menu application of \cref{thm:p17-maximal-choice}.
\end{hbmethod}

\begin{hbexample}{def:p17-vnm-worked-choice}{The same certain-prize ranking can select different risks}
Consider hypothetical consequences \(\ell,c,h\), ordered from worse to
better, and the following two lotteries:
\[
                 p=\tfrac12\delta_\ell+\tfrac12\delta_h,
                 \qquad q=\delta_c.
\]
Suppose the preference calibration gives
\(u(\ell)=0\), \(u(c)=2/5\), and \(u(h)=1\). Then
\(U_u(p)=1/2>2/5=U_u(q)\), so \(p\) is the selected member of
the menu \(\{p,q\}\). If instead the calibrated middle value is
\(3/5\), the same certain-prize ordering yields \(U_u(q)=3/5>1/2\),
and \(q\) is selected. These are two different lottery preferences.

The strictly increasing transformation \(f(t)=\sqrt t\) on \([0,1]\) makes this
distinction concrete. Applied to the first set of prize values it gives
expected scores \(1/2\) for \(p\) and \(\sqrt{2/5}>1/2\) for \(q\),
reversing their comparison. Applied to the completed lottery scores it gives
\(\sqrt{1/2}>\sqrt{2/5}\), preserving the original comparison.
The two operations are different.

For the parameterized lottery
\(p_\alpha=\alpha\delta_h+(1-\alpha)\delta_\ell\), the original
representation gives \(U_u(p_\alpha)=\alpha\). Thus
\(p_\alpha\succ q\) exactly when \(\alpha>2/5\), indifference holds
at \(2/5\), and \(q\succ p_\alpha\) below it. A positive affine
rescaling, for example \(v=5u-2\), preserves that threshold.
These stipulated numbers illustrate a decision rule; they are not measured
utilities or empirical findings.
\end{hbexample}

\begin{hbexample}{def:p17-common-consequence-example}{A common consequence constrains two choices}
Use the four lotteries from the source's common-consequence example,
with prize labels denoting monetary amounts
\cite[p.~101]{Rubinstein2015LectureNotes}:
\[
\begin{aligned}
 L_1&=\tfrac14\delta_{3000}+\tfrac34\delta_0,&
 L_2&=\tfrac15\delta_{4000}+\tfrac45\delta_0,\\
 L_3&=\delta_{3000},&
 L_4&=\tfrac45\delta_{4000}+\tfrac15\delta_0.
\end{aligned}
\]
One has \(L_1=\tfrac14L_3+\tfrac34\delta_0\) and
\(L_2=\tfrac14L_4+\tfrac34\delta_0\). Independence therefore requires
the same weak, strict or indifferent comparison in the two pairs.
For any consequence utility \(u\), direct finite summation gives
\[
                  U_u(L_1)-U_u(L_2)
                    =\tfrac14\bigl(U_u(L_3)-U_u(L_4)\bigr).
\]
The hypothetical paired choices \(L_2\succ L_1\) and \(L_3\succ L_4\)
cannot be represented by expected utility on this prize domain. No
assumption that utility equals money is needed for that conclusion.
The source also discusses observed responses; this example establishes
the mathematical incompatibility of the stated choices and does not
infer their population frequency or cause.
\end{hbexample}

\section{Why the assumptions matter}
\label{sec:p17-expected-utility-failures}

\begin{hbdefinition}{def:p17-vnm-counterpreferences}{Two alternative lottery preferences}
For a two-prize set \(Z=\{a,b\}\), define the score
\[
                       V(p)\defeq\max\{p(a),p(b)\},
\]
and prefer the lottery with the greater score, allowing ties. This is
the source's largest-probability criterion
\cite[pp.~95 and 99]{Rubinstein2015LectureNotes}.

For a three-prize set \(Z=\{a,b,c\}\), define a lexicographic preference:
two distinct vectors are compared at their first unequal coordinate in
the order \(a,b,c\); the vector with the larger coordinate is preferred.
Equal vectors are indifferent. This is the source's lexicographic example
\cite[pp.~96 and 99]{Rubinstein2015LectureNotes}.
\end{hbdefinition}

\begin{hbtheorem}{thm:p17-vnm-counterpreferences}{proved=annex:p17.vnm-counterpreferences}{Independence and continuity each exclude a different failure}
The largest-probability preference is complete, transitive and continuous,
but violates independence. The lexicographic preference is complete and
transitive and satisfies independence, but violates continuity.
Neither is represented by an expected-utility function on its stated
lottery domain. The failures have the following explicit witnesses.

For the first preference, \(\delta_a\sim\delta_b\), while
\[
 \tfrac12\delta_a+\tfrac12\delta_a
       \succ\tfrac12\delta_b+\tfrac12\delta_a.
\]
For the second, let
\[
 p=(1/2,1/2,0),\quad q=(1/2,0,1/2),\quad
 q_\varepsilon=(1/2+\varepsilon,0,1/2-\varepsilon),
                    \quad0<\varepsilon<1/2.
\]
Then \(p\succ q\), whereas \(q_\varepsilon\succ p\) for every such
\(\varepsilon\), and \(d(q_\varepsilon,q)\) tends to zero.
\end{hbtheorem}

\begin{hbwarning}{def:p17-vnm-interpretation}{What the representation warrants}
The theorem characterizes the preferences satisfying its stated axioms
on distributions of a fixed finite prize set. It does not prove that
people satisfy those axioms or that every decision should ignore the
random procedure, framing, or relevant state of the world. The source's
discussion of procedural fairness makes this limitation explicit
\cite[p.~102]{Rubinstein2015LectureNotes}.

The resulting scale is not a measure of interpersonal welfare. It also
does not provide subjective probabilities, risk-aversion comparisons on
a continuum of wealth, a theory of ambiguity, or the value of sequential
information. Those require their own domains, assumptions and proofs.
Rubinstein's separate betting model begins at pp.~102--103 and his
risk-aversion lecture at p.~108; their conclusions are not imported as
consequences of the present finite-prize theorem.
\end{hbwarning}
